\section{总结与延伸}

\subsection{全讲要点汇总}
\begin{table}[H]
\centering
\begin{tabular}{p{2.5cm}p{3.7cm}p{3.8cm}p{2.6cm}}
\toprule
\textbf{能力模块} & \textbf{课程核心判断} & \textbf{工程动作} & \textbf{主要风险} \\
\midrule
Reasoning & 为了更好行动，不是为了更长思维链 & 引入反思、重规划与动作校验 & Token 成本与幻觉放大 \\
Memory & 分层记忆是长任务稳定性核心 & 设计读写压缩淘汰策略 & 记忆污染与检索错配 \\
Planning & 动态规划优于静态清单 & 监控偏差并触发重规划 & 过度规划导致效率下降 \\
Multi-agent & 条件性增益，不是默认增益 & 先定协议，再增角色 & 通信开销与协作漂移 \\
Evaluation & 必须看全链路过程指标 & 建立 tracing 与统一 schema & 只看终局分数导致误判 \\
\bottomrule
\end{tabular}
\caption{Reasoning、Memory、Planning 的系统化落地框架}
\end{table}

\begin{importantbox}{一句话总结}
Language Agent 的下一阶段竞争，不是谁能做出更炫的 demo，而是谁能把推理、记忆和规划做成可验证、可运维、可扩展的工程系统。
\end{importantbox}

\subsection{延伸阅读}
\begin{itemize}
    \item Yao et al., \textit{ReAct: Synergizing Reasoning and Acting in Language Models}
    \item Shinn et al., \textit{Reflexion: Language Agents with Verbal Reinforcement Learning}
    \item Park et al., \textit{Generative Agents: Interactive Simulacra of Human Behavior}
    \item Wang et al., \textit{Voyager: An Open-Ended Embodied Agent with LLMs}
    \item OpenAI / Anthropic 关于 Agent tool-use 与 evaluation 的官方技术博客
\end{itemize}

\subsection{进一步实践建议}
\begin{itemize}
    \item 在单体 Agent 上先完成 ``可观测 + 可回滚'' 的最小闭环，再引入多智能体；
    \item 先建设任务级验证器，再扩大 reasoning budget；
    \item 将记忆操作显式 API 化，避免隐式上下文堆积；
    \item 对生产场景建立人工接管与降级策略，优先保障可控性。
\end{itemize}

\end{document}
